% from ega4-2-fr.tex
\oldpage[IV-2]{160}

\item[(ii)] Comme les $U_{C_n}(\mathcal F)$ sont ouverts pour tout $n$ par
(i), les $Z_n=X-U_{C_n}(\mathcal F)$ sont fermés dans $X$; en outre il est
clair que $Z_{n+1}\subset Z_n$, et comme $\dim(\mathcal F_x)$ est finie
pour tout $x\in X$ et
$\operatorname{coprof}(\mathcal F_x)\leq\dim(\mathcal F_x)$,
l'intersection des $Z_n$ est vide; comme on peut se borner au cas où $X$
est affine, donc quasi-compact, on peut supposer qu'il existe un $m$ tel
que $Z_m=\varnothing$. Or, il résulte de (5.7.4) que la relation
$x\in U_{S_k}(\mathcal F)$ équivaut à l'ensemble des relations
\begin{equation}
\operatorname{codim}_x(Z_n,S)>n+k
\tag{6.11.2.2}\label{IV.2.s6.11.2.2-fr}
\end{equation}
pour tout $n\geq0$; mais pour $n\geq m$ cette relation est automatiquement
vérifiée, donc on n'a en fait à considérer les relations (6.11.2.2) que
pour $0\leq n<m$. Or (0, 14.2.6) l'ensemble $V_{n,k}$ des $x$ vérifiant
(6.11.2.2) est ouvert, et $U_{S_k}(\mathcal F)$, intersection des $V_{n,k}$
pour $0\leq n<m$, est aussi ouvert. C.Q.F.D.
\end{enumerate}
\end{proof}

Rappelons que l'hypothèse que $X$ est localement immersible dans un schéma
régulier est toujours remplie lorsque $X$ est un préschéma localement de
type fini sur un corps $k$ (5.8.3).

\begin{corollary}[6.11.3]
\label{IV.2.s6.11.3-fr}
Sous les hypothèses de (6.11.2), l'ensemble $\operatorname{CM}(\mathcal F)$
des points $x\in X$ tels que $\mathcal F_x$ soit un module de
Cohen--Macaulay est ouvert dans $X$.
\end{corollary}

\begin{proof}
En effet, c'est l'ensemble $U_{C_0}(\mathcal F)$.
\end{proof}

\begin{remarks}[6.11.4]
\label{IV.2.s6.11.4-fr}
\begin{enumerate}
\item[(i)] Le raisonnement de (6.11.2, (ii)) prouve que (sans hypothèse sur
$X$) lorsque $U_{C_n}(\mathcal F)$ est ouvert pour tout entier $n$, alors
$U_{S_n}(\mathcal F)$ est ouvert pour tout entier $n$.
\item[(ii)] On a
$\operatorname{CM}(\mathcal F)=\bigcap_nU_{S_n}(\mathcal F)$. Si tout point
$x\in X$ admet un voisinage ouvert $V$ qui soit de dimension finie, la
suite des intersections $V\cap U_{S_n}(\mathcal F)$ est stationnaire
puisqu'il existe un $m$ tel que $\dim(\mathcal F_x)\leq m$ pour tout
$x\in V$; par suite, si les $U_{S_n}(\mathcal F)$ sont ouverts pour tout
entier $n$, $\operatorname{CM}(\mathcal F)$ est alors ouvert dans $X$.
\item[(iii)] On écrira $U_{S_n}(X)$, $U_{C_n}(X)$,
$\operatorname{CM}(X)$ au lieu de $U_{S_n}(\mathcal O_X)$,
$U_{C_n}(\mathcal O_X)$, $\operatorname{CM}(\mathcal O_X)$.
\end{enumerate}
\end{remarks}

\begin{proposition}[6.11.5]
\label{IV.2.s6.11.5-fr}
Soient $A$ un anneau local noethérien, $M$ un $A$-module de type fini. Pour
tout idéal premier $\mathfrak p$ de $A$, on a
\begin{equation}
\operatorname{coprof}_{A_{\mathfrak p}}(M_{\mathfrak p})\leq
\operatorname{coprof}_A(M).
\tag{6.11.5.1}\label{IV.2.s6.11.5.1-fr}
\end{equation}
\end{proposition}

\begin{proof}
On peut se borner au cas où $M_{\mathfrak p}\ne0$. Comme $\widehat A$ est
un $A$-module fidèlement plat (0$_{\rm I}$, 7.3.5), il existe un idéal
premier $\mathfrak q$ de $\widehat A$ au-dessus de $\mathfrak p$
(0$_{\rm I}$, 6.5.1); comme
$\widehat M_{\mathfrak q}=(M\otimes_A\widehat A)_{\mathfrak q}
=M\otimes_A\widehat A_{\mathfrak q}$ et que
$\widehat A_{\mathfrak q}$ est un $\widehat A$-module plat, donc un
$A_{\mathfrak p}$-module plat (0$_{\rm I}$, 6.2.1), il résulte de (6.3.2),
appliqué à l'homomorphisme local
$A_{\mathfrak p}\to\widehat A_{\mathfrak q}$, à $M_{\mathfrak p}$ et
$\widehat M_{\mathfrak q}$, que l'on a
$\operatorname{coprof}_{A_{\mathfrak p}}(M_{\mathfrak p})\leq
\operatorname{coprof}_{\widehat A_{\mathfrak q}}
(\widehat M_{\mathfrak q})$. D'autre part,
$X=\operatorname{Spec}(\widehat A)$ est isomorphe à un sous-schéma d'un
schéma régulier en vertu du théorème de Cohen (0, 19.8.8, (i)). On déduit
donc de (6.11.2) que l'on a
$\operatorname{coprof}_{\widehat A_{\mathfrak q}}
(\widehat M_{\mathfrak q})\leq
\operatorname{coprof}_{\widehat A}(\widehat M)$; enfin, on sait que
$\operatorname{coprof}_{\widehat A}(\widehat M)=
\operatorname{coprof}_A(M)$ (0, 16.4.10), ce qui achève de démontrer
(6.11.5).
\end{proof}

Cette proposition justifie la définition de la coprofondeur d'un
$A$-module lorsque $A$ n'est pas un anneau local, donnée dans (5.7.12).

\begin{proposition}[6.11.6]
\label{IV.2.s6.11.6-fr}
Soient $X$ un préschéma localement noethérien, $\mathcal F$ un
$\mathcal O_X$-Module cohérent, $n$ un entier $>0$. Supposons que pour tout
sous-préschéma fermé intègre $Y$ de $X$, il existe une partie ouverte non
vide $W$ de $Y$ telle que le sous-préschéma de $Y$ induit sur l'ouvert $W$
vérifie $(S_n)$. Alors l'ensemble $U_{S_n}(\mathcal F)$ est ouvert dans
$X$.
\end{proposition}

\begin{proof}
La question étant locale sur $X$, on peut se borner au cas où $X$ est
noethérien. Nous allons raisonner par récurrence sur $n$: pour $n=1$,
l'ensemble $U_{S_1}(\mathcal F)$ est ouvert dans $X$, car l'ensemble des
$x\in X$ où $\mathcal F$ ne vérifie pas $(S_1)$ est l'ensemble des $x$ tels
que $\mathcal F_x$ admette des idéaux premiers associés immergés (5.7.5);
si $(Z_j)$ est la famille finie des cycles premiers associés à $\mathcal F$
qui sont immergés, on a par suite
$U_{S_1}(\mathcal F)=X-\bigcup_jZ_j$, d'où notre assertion puisque les
$Z_j$ sont fermés. Nous supposerons donc désormais que $n>1$. En second
lieu, on peut se borner au cas où
$\operatorname{Supp}(\mathcal F)=X$, car il y a un sous-préschéma fermé
$T$ de $X$ ayant $\operatorname{Supp}(\mathcal F)$ pour espace sous-jacent,
et un $\mathcal O_T$-Module cohérent $\mathcal G$ tels que
$\mathcal F=j_*(\mathcal G)$, $j:T\to X$ étant l'injection canonique
(Err$_{\rm III}$, 30); comme il revient au même de prouver que $\mathcal F$
ou $\mathcal G$ vérifie $(S_n)$ en un point
$x\in\operatorname{Supp}(\mathcal F)$, on peut se borner à considérer le
cas où $T=X$. Notons enfin que par définition (5.7.2),
$U_{S_n}(\mathcal F)$ est stable par générisation. Nous allons utiliser le
lemme suivant:

\begin{lemma}[6.11.6.1]
\label{IV.2.s6.11.6.1-fr}
Soient $X$ un espace noethérien dans lequel toute partie fermée irréductible
admet un point générique, $E$ une partie de $X$. Pour que $E$ soit ouverte
dans $X$, il faut et il suffit que $E$ soit stable par générisation, et que,
pour toute partie ouverte $V$ de $X$ et toute partie irréductible $Y$ fermée
dans $V$, telles que $V-Y\subset E$ et que le point générique de $Y$
appartienne à $E$, $E\cap Y$ contienne une partie ouverte non vide de $Y$.
\end{lemma}

\begin{proof}
On observera que ce critère entraîne celui de (0$_{\rm III}$, 9.2.6)
lorsque toute partie fermée irréductible de $X$ admet un point générique.
Il n'y a évidemment à démontrer que la suffisance des conditions énoncées.

Considérons l'intérieur $U$ de $E$; l'ensemble fermé $X-U$ est réunion de
ses composantes irréductibles, qui sont en nombre fini et fermées dans $X$.
Si on avait $E\ne U$, l'hypothèse que $E$ est stable par générisation
entraînerait que le point générique $z$ d'une des composantes irréductibles
de $X-U$ appartiendrait à $E$. Or, $z$ n'appartient qu'à une seule des
composantes irréductibles de $X-U$; si $T$ est la réunion des autres
composantes irréductibles de $X-U$, $V=X-T$ est ouvert dans $X$, réunion de
$U$ et de l'ensemble $Y=\overline{\{z\}}\cap V$ fermé dans $V$ et
irréductible. Par hypothèse $E\cap Y$ contient une partie $W$ ouverte dans
$Y$; on en conclut que $U\cup W$ est ouvert dans $V$, donc dans $X$, ce qui
est absurde puisque $U$ est supposé être l'intérieur de $E$.
\end{proof}

